\section{Member 6: shared pipeline and the analytic model}
\label{sec:m6}
\textbf{Code:} \code{data.py}, \code{decision.py}, \code{metrics.py}, \code{config.py},
\code{plotting.py}, \code{run\_all.py}, \code{analytic.py}, \code{experiments/t6\_analytic.py}.

\paragraph{Data and splits.}
\emph{Spambase}~\cite{spambase} has $4{,}601$ emails, $1{,}813$ ($39.4\%$) of them spam, described by
57 word, character and capital-letter statistics. The features are heavily skewed and $77\%$ of all
values are exactly zero. \code{data.py} downloads and checksums the data, makes stratified
$60/20/20$ train/validation/test splits ($2{,}761/921/919$ emails), and fits the standardisation
(optionally after $\log(1+x)$) on the training part only. Every experiment repeats over the same 20
seeds and reports the mean and its standard error. Models are fitted on training data, every
choice (penalty, shrinkage, temperature) is made on validation data, and the test part is used only
to report. \code{data.py} also draws synthetic Gaussian data, whose exact posterior is known, and
resamples a test set within each class to a new spam rate, which Member 1 needs in
Section~\ref{sec:m1}.

\paragraph{Decision rule and metrics.}
\code{decision.py} turns any loss matrix into the thresholds of Section~\ref{sec:intro} and applies
them, either to probabilities or directly to log-odds (which avoids rounding a confident score to
exactly $0$ or $1$). A unit test compares it with a brute-force search over the three actions on
$10{,}000$ random posteriors and loss matrices. \code{metrics.py} computes the realized loss, accuracy,
inspection rate and confusion matrices with an inspect column. \code{run\_all.py} regenerates every
figure and table in this report in under a minute, and 101 unit tests cover the shared code and
the members' modules.

\paragraph{An analytic check: accuracy against cost.}
To see the trade-offs without any fitting, \code{analytic.py} uses a one-dimensional model:
real emails have a score $x\sim N(-1,1)$, spam has $x\sim N(+1,1)$, and $40\%$ of emails are spam.
Here the accuracy, cost and inspection rate of any threshold rule have closed forms.
Table~\ref{tab:t6} compares four rules.

\begin{table}[!htbp]
\centering\small
\input{\resultsdir/t6_rule_comparison.tex}
\caption{Exact performance in the one-dimensional model ($\cfp=10$, $\cfn=1$). FPR: share of real
emails sent to spam; FNR: share of spam delivered.}
\label{tab:t6}
\end{table}

\begin{figure}[!htbp]
\centering
\includegraphics[width=0.92\linewidth]{t6_accuracy_vs_cost.pdf}
\caption{One-dimensional model: accuracy and expected cost of the rule ``spam iff $\post\ge t$''.
Accuracy peaks at $t=0.5$, cost is lowest at $\tstar=10/11$.}
\label{fig:t6-acc}
\end{figure}

The accuracy rule is the most accurate automatic rule ($84.6\%$) but costs $0.772$ per email,
more than delivering everything ($0.400$). Figure~\ref{fig:t6-acc} shows why: between $0.5$ and
$10/11$, every step that raises accuracy also raises the cost. Moving the cut to $\tstar=10/11$ lowers accuracy to
$73.9\%$ and the cost to $0.311$, because it almost never files a real email as spam (FPR $0.9\%$).
Allowing inspection at $\ci=0.3$ cuts the cost to $0.179$ while $41.6\%$ of emails are inspected.
Figure~\ref{fig:t6} shows how this depends on the price of inspection. The best achievable cost rises
with $\ci$, and its slope at any price equals the share of emails inspected at that price. Inspection
stops entirely at $\ci=\cstar=10/11$, whatever the class separation.

\begin{figure}[!htbp]
\centering
\includegraphics[width=0.92\linewidth]{t6_risk_reject_vs_cinsp.pdf}
\caption{One-dimensional model: lowest achievable cost (left) and share of emails inspected (right)
against the inspection price $\ci$, for three class separations $\mu$. Dotted: the best cost
without inspection. Dashed: $\cstar=10/11$, where inspection stops.}
\label{fig:t6}
\end{figure}
